\chapter{Advanced Rotorcraft Transmission (A.R.T.) program}
\section{Introduction}
	The Advanced Rotorcraft Transmission (A.R.T.) program was created to develop and demonstrate lightweight, 
	quiet, durable drive train systems for the next generation of rotorcraft. Sikorsky is one of the participants, 
	and this program allows them to evaluate key emerging material and component technologies and novel design 
	concepts for advancing the technology for future rotorcraft transmission.  

	One of the specific objectives that is relevant to the project is “Reduce transmission weight by 25 percent 
	relative to SOA — the most recent stage in the development of a product, incorporating the newest technology, 
	ideas, and features — trends (currently in the range of 0.40 lb./hp)”. A.R.T. is designed to provide a high 
	return on weight improvement by developing novel concepts involving high risk factors. 

	The A.R.T program requires all participants to select a vehicle fitting into one of the two distinct 
	next-generation aircraft classes and to demonstrate their proposed advanced-concept drive trains. The first one 
	is Future air attack vehicle (FAAV): a 10,000- to 20,000-pound aircraft capable of undertaking tactical support 
	and air-to-air missions, and second is advanced cardo aircraft (ACA): a 60,000- to 80,000-pound aircraft 
	capable of heavy-lift field support operations. 

	To evaluate the effect of advanced-technology improvements on aircraft size and performance, computer 
	programs were initially used to develop the reference aircraft design. In the early stage of this program, 
	a few key parameters were used to develop the drive system design. As the design progressed, detailed rotor 
	and engine performance data were entered to improve the precision of the calculations. The program was used to 
	discover broad trends early in the design process, and it was also used to determine the effect of relatively 
	small modifications once the design was more developed. 

	Included in the chapter is also the trade studies selection process. As during the early development of the 
	A.R.T. drive system, comparisons and trade studies were conducted among novel transmission configurations with 
	high payoff potential. 

\section{Program Objectives}
	While standard industrial gearboxes prioritize manufacturing cost and traditional safety factors, 
	aerospace drivetrains demand ultra-lightweight, quiet, and highly durable systems. The Advanced Rotorcraft 
	Transmission (A.R.T.) program—a joint initiative funded by the U.S. Army and NASA—was established to drive 
	drivetrain technology beyond historical baseline limits. The program set strict performance targets to improve 
	on contemporary baseline aircraft metrics: 
		\begin{itemize}
			\item \textbf{Weight Reduction:} Lower total transmission weight by \textbf{25\%} relative to prevailing baseline configurations, which typically trended in the range of 0.40 lb/hp.
			\item \textbf{Acoustic Attenuation:} Suppress operational noise levels by \textbf{10 dB} at the transmission source, reducing acoustic fatigue on components and passengers.
			\item \textbf{Operational Reliability:} Extend the Mean Time Between Removals (MTBR) to a minimum of \textbf{5,000 hours} to lower lifecycle costs and enhance mission readiness.
		\end{itemize}
	Drivetrain weight reductions have direct, compound benefits for overall aircraft performance. Retrofitting a 
	baseline helicopter with an advanced A.R.T. transmission allows saved empty weight to be converted into useful 
	load capacity without increasing gross takeoff weight. This capacity can be allocated directly to cargo payload 
	or fuel volume. Increasing fuel capacity extends the aircraft's operational radius of action, allowing for 
	longer flight profiles and enhanced mission flexibility. 

\section{Advanced Structural Drivetrain Configurations}
	To meet these targets, contractors evaluated and validated several advanced architectural configurations: 
		\begin{enumerate}
			\item \textbf{Self-Aligning Bearingless Planetary (SABP) Configuration}\par
			This layout replaces conventional, heavy planet roller bearings and heavy carrier assemblies with an 
			internal design utilizing flexible structural rings. By allowing the planet gears to float and self-align 
			within the mesh, the system automatically balances tooth loads across all paths. This architecture eliminates 
			high-wear bearings, lowers structural weight, and improves system reliability. 
			\item \textbf{Split-Torque Configuration}\par
			The split-torque design divides the incoming power from a single engine input into multiple low-power, 
			parallel gear paths. These separate branches recombine a single output gear. Splitting the torque allows for 
			smaller gear faces and narrower tooth profiles, which reduces total system weight and provides redundant load paths. 
			\item \textbf{Conventional Single-Stage Planetary Configuration (Aerospace Optimized)}\par
			This approach refines the standard planetary layout using advanced high-strength steel alloys and optimized 
			micro-geometry. It maintains central sun gear, external stationary ring gear, and a planet carrier. 
			This baseline configuration was used as a benchmark to measure the weight and noise advantages of alternative architectures. 
		\end{enumerate}

\section{Program Development-Sikorsky}
	The major objective of the A.R.T. contract was the development of a high reduction ratio, torque splitting output stage. 
	For minimum weight, the output stage of the split path design has a high aspect ratio of face width to diameter.  

	“Preliminary designs were conducted of split path and split torque transmissions to evaluate the weight, reliability, 
	and noise and compare the improvements to the program goals. A split torque gearbox with a high reduction ratio double 
	helical output gear was determined to be 23% lighter, greater than 10dB quieter, and almost four times more reliable than 
	the baseline two stage planetary design. The improvements were attributed to extensive use of composites, spring clutches 
	located at the transmission input, advanced high hot hardness gear steels, the split path configuration itself, 
	high reduction ratio, double helical gearing on the output stage, elastomeric load sharing devices, and elimination of 
	accessory drives.” 

	Baseline gearbox weight was determined using statistical weight equations that utilize preliminary weight parameters 
	such as design torque and reduction ratio. Layout weight calculations were conducted on the advanced split path and split 
	torque designs to obtain the individual component weights and to enable identification of bulk items which the layout 
	does not include such as hardware, protective coatings, sealants, bosses on castings, lubrication system components, 
	oils, etc. Once identified, the bulk item weights are estimated using values obtained from systems of similar design. 
	These layout calculations are compared to the analytical weight and differences resolved. 

	Since one engine inoperative (OEI) power is felt all the way to the final stage bull gear mesh in the split path 
	type of arrangement, the gearbox is essentially sized to OEI conditions. 

	A cross-sectional layout of the ACA baseline main gearbox having the configuration of the CH53E is 
	depicted in \autoref{fig:gearbox3} which also shows the rpms, reduction ratio's, numbers of teeth, and pitch diameters for the main 
	meshes. This gearbox has been sized to the speed and power requirements of the ACA.

		\begin{figure}[h]
			\centering
			\includegraphics[width=0.9\textwidth]{gearbox/gearbox3}
			\caption{ACA Baseline Transmission}
			\label{fig:gearbox3}
		\end{figure}
		\FloatBarrier

\section{Approach and Methodology}
	The A.R.T. weight savings results in additional reductions when the ACA aircraft is downsized for the same payload and 
	performance. Weight savings from the transmission alone is 2,508 pounds. If the aircraft size is not changed, the payload 
	will increase by the transmission weight savings. If the payload is kept constant, the aircraft can be downsized.

	Note that the weight empty is reduced by 2,508 pounds for the baseline w/A.R.T. since this is the savings for the A.R.T. 
	transmission system alone. There is however an additional 1,201 pounds savings by downsizing the aircraft for the same 
	payload thus the ART has actually saved 3,709 pounds in the aircraft.  When resizing the aircraft for the same payload, 
	certain aircraft design parameters were kept constant. The engines, three GE-38's, were fixed for all designs. The cabin 
	size was fixed. The payload/range was kept at 25,000 pounds with a radius of action of 500 kilometers. The takeoff 
	criterion of 200 feet per minute vertical rate of climb at 95% internal rated power, 4,000 feet, 95°F day was kept 
	constant. With the disk loading maintained at approximately I0 pounds per square foot, the main rotor diameter was 
	sized to match the power required with the power available at the takeoff condition.  

		\begin{figure}[h]
			\centering
			\includegraphics[width=0.9\textwidth]{gearbox/gearbox4}
			\caption{Weight Saving by Downsizing Aircraft for Same Payload}
			\label{fig:gearbox4}
		\end{figure}
		\FloatBarrier

	All aircraft design analysis was accomplished with the helicopter design model (HDM). The HDM is a computer 
	program which uses basic input data for mission requirements, design criteria, configuration, engine, and technology. 
	The program then calculates dimensions, weights, powers, performance parameters, and costs. The HDM uses experience 
	factors for calculating many of the outputs. For example, weight trending equations are provided for the main gearbox, 
	main rotor shaft, main rotor blades, intermediate gearbox, tail gearbox, and tail drive shafting. These equations use 
	parameters such as horsepower, rpm, and reduction ratio with coefficients derived from previous designs to calculate weights. 

		\begin{figure}[h]
			\centering
			\includegraphics[width=0.9\textwidth]{gearbox/gearbox5}
			\caption{Helicopter Design Model}
			\label{fig:gearbox5}
		\end{figure}
		\FloatBarrier

\newpage
\section{Trade Studies}
		Four preliminary designs were completed in the preliminary design phase of the program. Designs of split path and split torque 
		gearboxes were completed for a high output bull gear reduction ratio of 10.85:1 (141/13) and a lower ratio of 8.15:1 (106/13) for a 
		total of two ratios of split torque and two ratios of split path designs. 

		When they entered the detail design phase, more realistic calculations revealed that the high ratio configuration (10.85:1) produced 
		excessive stresses in the bull pinion shaft (the shaft that drives the big “bull” gear). To stay within safe stress limits, the maximum 
		ratio was retuned down to 9.77:1 (which is the ratio 127/13). 

		Split-Torque uses a planetary/“differential” arrangement to split the engine's torque before it goes to the main gearbox. 
		A differential type planetary gear set splits the torque that comes from the two engines. Split-Path has two engine drives which go 
		through separate paths that meet only at the final “bull” gear that delivers power to the rotor. This means that it runs the two 
		engine drives on separate mechanical paths that only meet in the very last gear (the bull gear). 

		The trade-off analysis ultimately selected the split-path gearbox powered by a lower reduction ratio (106/13) because it is the 
		only architecture that simultaneously delivers the lightest weight, the highest reliability (MTBR), room-for-growth in manufacturing, 
		and compliance with the low-noise goal. The split-torque design, despite its inherent torque balance, fell short due to extra gear 
		stages, heavier bearing loads, and reduced redundancy; all of which would have increased weight and lowered MTBR for the ACA's 
		three-engine configuration. 